\documentclass{Academic}

\begin{document}
%Easy customisation of title page
%TC:ignore
\myabstract{\small
\noindent\textbf{Resumen}
%Abstract below

La enfermedad renal crónica (ERC) es una pérdida progresiva de la función renal que afectó a más de 800 millones de personas en 2017, 
con frecuencia subdiagnosticada por presentar síntomas en etapas avanzadas. Su detección temprana podría reducir significativamente su 
carga clínica y económica. El siguiente estudio empleó un enfoque cuantitativo no experimental basado en redes bayesianas gaussianas (GBN) 
para modelar las relaciones de dependencia probabilística entre biomarcadores clínicos asociados a la ERC, utilizando la base de datos 
ckd.csv. Se construyeron tres grafos acíclicos dirigidos (DAGs) con distintos subconjuntos de variables, definidos con apoyo de estudiantes 
y profesores de medicina, ajustados con bn.fit() en R mediante el paquete bnlearn y evaluados utilizando los criterios BIC y AIC. La DAG 2 
presentó el mejor desempeño global y fue seleccionada para realizar las consultas probabilísticas propuestas por especialistas. Los 
resultados mostraron probabilidades de 62.77 % de creatinina sérica elevada dado un nivel alto de sodio, 62.84 % al considerar 
simultáneamente sodio y potasio elevados, 43.13 % de hiperpotasemia ante creatinina y urea elevadas, 80.96 % de hemoglobina menor a 12 g/dL 
ante una creatinina de 3.0 mg/dL y un conteo de eritrocitos de 3.5 millones/\(\mu\)L, y 100 % de creatinina elevada ante presión arterial 
superior a 140 dentro del modelo ajustado. Asimismo, el modelo no paramétrico mejoró los valores de BIC y AIC de la DAG 2 respecto al modelo 
lineal, aunque las consultas probabilísticas se realizaron sobre este último. Finalmente, se identificó como limitación la imposibilidad de 
incorporar directamente variables categóricas, como el diagnóstico de ERC, en una red exclusivamente gaussiana, proponiéndose como extensión 
futura el uso de redes bayesianas gaussianas condicionales o redes discretas.
}
\renewcommand{\myTitle}{Uso de Redes Bayesianas Gaussianas para comprender 
el efecto de biomarcadores en la Enfermedad renal crónica: Un análisis probabilístico}
\renewcommand{\MyAuthors}{%
    Jessica Milagros Cordero Rivera$^{1}$ \quad
    Diego Rodrigo González Ramírez$^{1}$\\[0.2em]
    Jose Ángel Lechuga Carrasco$^{1}$ \quad
    Julio Lizárraga Beltrán$^{1}$
}

\renewcommand{\MyAffiliations}{%
    $^{1}$Escuela de Ingeniería y Ciencias
}

\renewcommand{\AuthorIDs}{%
    A01647083 \;|\; A01647198 \;|\; A00228198 \;|\; A01743628
}

\renewcommand{\ShortAuthors}{Cordero Rivera et al.}

\renewcommand{\Keywords}{enfermedad crónica renal, red bayesiana, red bayesiana gaussiana, keyword4, keyword5}
\maketitle
%\vspace{-1.9em}\noindent\rule{\textwidth}{1pt} %add this line if not using abstract
\onehalfspacing
%TC:endignore

\section{Introducción}


La enfermedad renal crónica (ERC) es la pérdida progresiva de la función de los riñones, 
afectando su función de eliminar residuos y exceso de agua del cuerpo. Al encontrarse
en una etapa tardía del ERC el paciente puede experimentar síntomas como, entre otros:

\begin{itemize}
  \item Piel anormalmente oscura o clara
  \item Dolor de huesos
  \item Problemas para concentrarse o pensar
  \item Fasciculaciones musculares o calambres
  \item Susceptibilidad a hematomas
  \item Sed excesiva
  \item Hinchazón en las manos y en los pies
\end{itemize}

A muchas personas no se les diagnostica la ERC sino hasta que han perdido 
gran parte de su función renal.\cite{definicion}. 
En 2017 la ERC afectó a más de 800 millones de personas y su tasa de mortalidad ha 
estado en aumento en este siglo. Además, debido a otras enfermedades cómo la diabetes mellitus 
y la obesidad, el riesgo de padecer ERC aumenta aún más y con ello, sus complicaciones \cite{Csaba}.

Es evidentemente claro que la identificación y diagnóstico de la ERC beneficiaría; para unas, 
La calidad de vida de millones de personas que podrían desarrollar esta enfermedad en un 
futuro; y para otras, evitar fallecer por esta enfermedad o complicaciones derivadas de ella. 
Los estudios epidemiológicos han concientizado sobre el problema de la ERC no 
diagnosticada y sugieren que su detección y tratamiento tempranos reducirán
la carga mundial que suponen los pacientes que necesitan diálisis \cite{Taal}.

Sin embargo,


\section{Metodología}

\subsection{Enfoque y tipo de estudio}
Se empleó un enfoque cuantitativo, no experimental, basado en el uso de GBN para encontrar 
las relaciones de dependencia probabilística entre los biomarcadores

\subsection{Fuente de datos}
Se nos otorgo una base de datos en un archivo .csv con el nombre ckd.csv, esta base de datos cuenta con las siguientes variables:
\begin{itemize}
    \item age: Age.
    \item bp: Blood pressure.
    \item sg: Specific gravity.
    \item al: Albumin.
    \item su: Sugar.
    \item rbc: Red blood cells.
    \item pc: Pus cell.
    \item pcc: Pus cell clumps.
    \item ba: Bacteria.
    \item bgr: Blood glucose random.
    \item bu: Blood urea.
    \item sc: Serum creatinine.
    \item sod: Sodium.
    \item pot: Potassium.
    \item hemo: Hemoglobin.
    \item pcv: Packed cell volume.
    \item wc: White blood cell count.
    \item rc: Red blood cell count.
    \item htn: Hypertension.
    \item dm: Diabetes mellitus.
    \item cad: Coronary artery disease.
    \item appet: Appetite.
    \item pe: Pedal edema.
    \item ane: Anemia.
    \item class: Diagnosis.
\end{itemize}
Adicionalemente a esto, entrevistamos a profesores y alumnos de la escuela de medicina del Instituto Tecnológico y de Estudios Superiores de Monterrey, que esten familiarizados con el ERC, para que nos ayudaran a generar 3 DAG`s que represente las relaciones de dependencia, ademas de esto, a las personas mas especialistas en el tema se les solicito que nos ayudaran a proponer 2 queries para poderlas resolver utilizando nuestro modelo.


\subsection{Herramientas}
Todos los análisis probabilisticos se realizaron utilizando entorno de desarrollo 
integrado \texttt{RStudio} (v2026.08.2), utilizando el lenguaje de programación 
estadístico \texttt{R} (v4.5.3) \cite{R} mediante el sistema de código abierto 
Quarto (v1.8.25), utilizando los paquetes \texttt{bnlearn} (v5.2.1)\cite{Bnlearn} 
y \texttt{tidyverse} (v2.0.0).

\subsection{Preprocesamiento}
Durante este proceso se redujo el tamaño de las DAG`s para que solo tuvieran las variables que se iban a utilizar para la DAG respectiva, asi mismo, se limpiaron los datos en todas las DAG`s utilizando la funcion na.omit en Rstudio para que no existiera ninguna fila con datos vacios. Para la primer DAG solo se seleccionaron las variables siguientes: \texttt{sod, pot, bp, sc, bu}. Para la segunda DAG se seleccionaron las siguientes: \texttt{sod, bp, sc}. Y para la tercer DAG se utilizaron: \texttt{bgr, su, sg, al, sc, bu, pot, sod, bp, rc, hemo, pcv}, para esta tercer DAG tambien se convirtieron todas las variables a tipo numericas.

\subsection{Construcción y selección del modelo global}
Para la construcción de las DAGs se utilizó el criterio de 3 estudiantes de la escuela de 
medicina de 5to semestre, 2 del Tecnológico de Monterrey campus Guadalajara y 1 de 
la Universidad Autónoma de Guadalajara. Cada red fue ajustada utilizando la funcion de 
bnlearn \texttt{bn.fit()} con los valores numéricos de la base de datos \texttt{ckd.csv}
y evaluada con los criterios de información BIC y AIC mediante \texttt{score()}.

\subsection{Estrategia analítica por pregunta de investigación}
Para cada pregunta se utilizo un DAG distinto tomando en cuenta lo que nos interesaba 
saber, se ajstaron con \texttt{bn.fit()} y se encontraron probabilidades condicionales 
con \texttt{cpquery()} (utilizando una \(n=10^{6}\) y verificado con ponderación por
verosimilitud, \texttt{method="lw"}).


\section{Resultados}




\section{Discuciones}


\section{Conclusiones}


%TC:ignore
%\clearpage %add new page for references
\singlespacing
\emergencystretch 3em
\hfuzz 1px
\printbibliography[heading=bibnumbered]

% \clearpage
% \begin{appendices}

% \section{Here go any appendices!}

% \end{appendices}

%TC:endignore
\end{document}